% !TeX encoding = utf-8
% !TeX program = pdflatex
%
% 数维杯 B 题论文主文件（骨架）。
% 编译：pdflatex（推荐）或 xelatex，需与 swmcmthesis.cls / swmcmthesis.cfg 同目录。
% 图片：优先从 ../figures/（仓库根 figures/）查找。
% 注意：正文模型/算法/结论由 A、B 撰写，本骨架只提供结构与符号表占位。

\documentclass{swmcmthesis}

% 图片搜索路径（根 figures/ 与 paper 内 figures/ 都找）
\graphicspath{{../figures/}{figures/}}

%%%%%%%%%%%%填写相关信息%%%%%%%%%%%%%%%%%%%%%%%%%%
\title{Rapid UAV Localization via Coordinated Mobile Vehicles under Occlusion and Uncertainty} % TODO 论文题目，可再改

\baominghao{SWBMCM2600001}           % TODO 改成自己的报名号

\zubie{B}                            % 选题 B

\abstract{%
% TODO 摘要最后写：须覆盖四问的方法与主要结果（P1 可观测区域与最优站位、P2(1) 精度指标、
% P2(2) 时间线与基线对比、P3 误差门槛、P4 的 η 与 t_last）。不要出现学校/导师/队员姓名。
Summary of the results goes here.
}

\keywords{UAV localization; bearing intersection; observable region; vehicle routing; rescheduling}

\date{2026}

\begin{document}

% 摘要页
\maketitle

% 目录
\tableofcontents
\newpage

% 正文
\setcounter{page}{1}

\section{Introduction}
% TODO C 写：背景（低空无人机、粗告警、双站方位交会）、研究目标、四问概述
\subsection{Problem Background}
% TODO
\subsection{Restatement of the Problem}
% TODO 用符号重述四问
\subsection{Our Work}
% TODO 方法与结论概览

\section{Assumptions and Notations}
\subsection{Assumptions}
% TODO A/B 提供模型假设（与 shared/constants.md 保持一致）
\begin{assumption}
% 示例：车辆匀速 3 m/s，忽略加减速与转向时间。
\end{assumption}

\subsection{Notations}
% 符号表（与 shared/constants.md 一致，待 A/B 复核）
\begin{table}[!h]
\centering
\caption{Main notations}
\begin{tabular}{ll}
\toprule[1.5pt]
Symbol & Description\\
\midrule[1pt]
$O$ & Dispatch center, $O=(50,50)$ m\\
$B_k$ & Building $k$, $k=1,\dots,5$\\
$V_j$ & Observation vehicle $j$, $j=1,\dots,5$\\
$T_i$ & UAV target $i$, $i=1,2,3$\\
$\theta$ & Alert azimuth, counterclockwise from due east ($^\circ$)\\
$R$ & Alert horizontal distance (m)\\
$H$ & Alert target altitude (m)\\
$r$ & Horizontal distance from camera to target (m)\\
$\phi$ & Line-of-sight elevation angle ($^\circ$)\\
$\gamma$ & Intersection angle of two bearing lines ($^\circ$)\\
$L$ & Baseline length between two stations (m; problem uses $B$)\\
$Q$ & Observation quality metric, $Q=1/r$\\
$s$ & Cumulative travel distance (m)\\
$\Delta\psi$ & Horizontal bearing deviation ($^\circ$)\\
$\sigma(s)$ & Localization error std, $\sigma(s)=a+bs$ (m)\\
$\eta$ & Completion ratio in evaluation window\\
$t_{\mathrm{last}}$ & Latest observation start time (s)\\
$t_{\mathrm{rec}}$ & Communication recovery time (s)\\
\bottomrule[1.5pt]
\end{tabular}
\end{table}

\section{Task 1: Observable Region and Optimal Station (P1)}
% TODO B 撰写：三类可观测性约束、可观测区域、Q=1/r 最优站位及俯仰最小半径、区域图

\section{Task 2: Two-Station Localization and Vehicle Assignment (P2)}
\subsection{Accuracy Model and Optimal Station Pair (P2(1))}
% TODO B 撰写：精度指标定义、基线/交会角/距离的关系、30 s 行驶约束、退化情形讨论
\subsection{Vehicle Reuse and Routing under Concurrent Alerts (P2(2))}
% TODO A 撰写：车辆指派+路径规划、min-max 目标、每车完整时间线、与最近指派启发式对比

\section{Task 3: Error Propagation on Horizontal Bearing Accuracy (P3)}
% TODO B 撰写：三误差源、Δψ 传播模型、P(|Δψ|≤5°)≥0.95 的最大行驶距离、σ(s)=a+bs 约束

\section{Task 4: Rescheduling under Disturbances (P4)}
% TODO A 撰写：新告警/单车故障/通信中断及组合、场景集设计、η 与 t_last、重调度 vs 保留原计划

\section{Sensitivity Analysis}
% TODO 三人：对交会角、误差参数、扰动参数等的敏感性

\section{Strengths and Weaknesses}
% TODO 三人

\section{Conclusion}
% TODO 三人

% 参考文献
\begin{thebibliography}{9}
% TODO 填入实际引用文献
\bibitem{ref1} Reference entry goes here.
\end{thebibliography}

% 附录：核心代码
\newpage
\section*{Appendix}
\addcontentsline{toc}{section}{Appendix}
% TODO 关键程序以 lstlisting 放入附录，例如：
% \begin{lstlisting}[language=python,caption={Observable region computation}]
% ...
% \end{lstlisting}

\end{document}
